\documentclass{amsart}
\usepackage{amsmath,amssymb,amsthm}

\title{The Sinc$^2$ Continuum Limit and Fej\'er Kernel Identity}
\author{Brayden Ross Sanders / 7Site LLC}
\date{2026}

\newtheorem{theorem}{Theorem}
\newtheorem{corollary}[theorem]{Corollary}
\newtheorem{definition}[theorem]{Definition}

\begin{document}
\maketitle

\begin{abstract}
We prove two results connecting the First-G resonance function $R(k,f)$ to
classical analysis. (1) The sinc$^2$ continuum limit: $R(k,f) \to \mathrm{sinc}^2(t)$
as $f \to \infty$ with $t = kf$ fixed, at rate $O(1/f^2)$.
(2) The Fej\'er identity: $k \cdot R(k,f) = F_k(f)$, the Fej\'er kernel.
\end{abstract}

\section{Resonance Function}

\begin{definition}[Resonance function]
For $k \geq 1$ and $f > 0$ with $f \notin \mathbb{Z}$:
\[
  R(k,f) = \frac{\sin^2(\pi k f)}{k^2 \sin^2(\pi f)}.
\]
\end{definition}

\section{Sinc$^2$ Continuum Limit}

\begin{theorem}[Sinc$^2$ limit]
Fix $t \in (0,1)$, $t \notin \mathbb{Z}$. As $f \to \infty$ with $k/f = t$ fixed:
\[
  R(k,f) \longrightarrow \mathrm{sinc}^2(t) = \frac{\sin^2(\pi t)}{(\pi t)^2},
  \quad \text{at rate } O(1/f^2).
\]
\end{theorem}

\begin{proof}
Write $f = k/t$, so $\varepsilon = \pi/f \to 0$.
\[
  R(k,f) = \frac{\sin^2(\pi t)}{k^2 \sin^2(\varepsilon)}.
\]
Since $\sin(\varepsilon) = \varepsilon - \varepsilon^3/6 + O(\varepsilon^5)$,
we have $\sin^2(\varepsilon) = \varepsilon^2(1 - \varepsilon^2/3 + O(\varepsilon^4))$,
so $k^2 \sin^2(\varepsilon) = (k\varepsilon)^2(1+O(\varepsilon^2)) = (\pi t)^2(1+O(1/f^2))$.
Therefore $R(k,f) = \sin^2(\pi t)/((\pi t)^2(1+O(1/f^2))) \to \mathrm{sinc}^2(t)$.
\end{proof}

\begin{corollary}
$\mathrm{sinc}^2(1/2) = 4/\pi^2$ exactly, since $\mathrm{sinc}(1/2) = \sin(\pi/2)/(\pi/2) = 2/\pi$.
\end{corollary}

\section{Fej\'er Kernel Identity}

\begin{theorem}[Fej\'er identity]
For all $k \geq 1$ and $f > 0$ with $f \notin \mathbb{Z}$:
\[
  k \cdot R(k,f) = F_k(f),
  \quad \text{where } F_k(f) = \frac{1}{k} \cdot \frac{\sin^2(\pi k f)}{\sin^2(\pi f)}
\]
is the Fej\'er kernel of order $k$.
\end{theorem}

\begin{proof}
$k \cdot R(k,f) = k \cdot \dfrac{\sin^2(\pi k f)}{k^2 \sin^2(\pi f)}
= \dfrac{\sin^2(\pi k f)}{k \sin^2(\pi f)} = F_k(f)$.
\end{proof}

\begin{corollary}[Gap boundary is a zero]
At the First-G event $(k,f) = (p, 1/p)$:
\[
  R(p, 1/p) = \frac{\sin^2(\pi p \cdot 1/p)}{p^2 \sin^2(\pi/p)}
             = \frac{\sin^2(\pi)}{p^2 \sin^2(\pi/p)} = 0,
\]
since $\sin(\pi) = 0$ exactly. The First-G boundary is a zero of $R$.
In the continuum limit, $t = kf = p \cdot (1/p) = 1$ maps to $\mathrm{sinc}^2(1) = 0$,
the first zero of the sinc$^2$ function.
\end{corollary}

\end{document}
